% Generated draft: requires mathematical and source review.
\documentclass[11pt]{article}
\usepackage[margin=25mm]{geometry}
\usepackage{fontspec}
\setmainfont{Latin Modern Roman}
\usepackage{amsmath,amssymb,mathtools,unicode-math}
\setmathfont{Latin Modern Math}
\usepackage{enumitem,needspace,xurl,hyperref}
\setlength{\parindent}{0pt}
\setlength{\parskip}{5pt}
\setlength{\emergencystretch}{3em}
\newcommand{\N}{\mathbb{N}}
\newcommand{\Z}{\mathbb{Z}}
\newcommand{\Q}{\mathbb{Q}}
\newcommand{\R}{\mathbb{R}}
\newcommand{\C}{\mathbb{C}}
\newcommand{\F}{\mathbb{F}}
\newcommand{\A}{\mathbb{A}}
\newcommand{\PP}{\mathbb{P}}
\newcommand{\E}{\mathbb{E}}
\newcommand{\eps}{\varepsilon}
\newcommand{\dd}{\,\mathrm{d}}
\begin{document}
\section*{UnsolvedMath ID: 30003839}
\textbf{Source problem code:} catalogue entry
\subsection{Definitions and mathematical statement}
Let \(D=\Z[1/2]=\{m/2^k:m\in\Z,\ k\in\Z_{\geq 0}\}\). Thompson's group \(F\) is the group, under composition, of increasing piecewise linear homeomorphisms \(f:[0,1]\to[0,1]\) with finitely many breakpoints, all breakpoints in \(D\cap(0,1)\), and every linear slope in \(\{2^n:n\in\Z\}\). In particular, \(f(0)=0\) and \(f(1)=1\).

For words \(u,v\) over a finite alphabet, their padded convolution pairs corresponding letters, padding the shorter word on the right with a new symbol \(\square\) until both words have the same length. A binary relation on words is synchronously regular if its set of padded convolutions is accepted by a finite automaton.

Use the following standard finite-state interpretations of the two structures in the question. A finitely generated group \(G\) is automatic if there exist a finite symmetric generating set \(S\) and a regular language \(L\subseteq S^*\) such that the evaluation map \(\pi:L\to G\) is bijective and, for every \(s\in S\), the relation
\[
R_s=\{(u,v)\in L\times L:\pi(u)s=\pi(v)\}
\]
is synchronously regular. Here words in \(L\) represent the products of their generator letters; the use of unique representatives is a standard equivalent convention for automatic groups.

The group \(G\) admits a graph-automatic structure if there exist a finite symmetric generating set \(S\), a finite alphabet \(A\), a regular language \(K\subseteq A^*\), and a bijection \(\nu:K\to G\) such that, for every \(s\in S\),
\[
E_s=\{(u,v)\in K\times K:\nu(u)s=\nu(v)\}
\]
is synchronously regular. The labels in \(K\) need not be words evaluating to their assigned group elements. Thus this is a finite-state presentation of the generator-labelled Cayley graph.

\textbf{Questions.} The supplied source asks: ``Why should Thompson's group \(F\) fail to be automatic, and does \(F\) admit a graph-automatic structure?'' Under the conventions above, the first question asks whether \(F\) has no automatic structure and, if that expected negative answer is correct, for a proof of nonexistence. Its phrasing does not supply a theorem that \(F\) is nonautomatic. The second question asks whether there exist \(S,A,K,\nu\) satisfying the graph-automatic conditions above for \(G=F\).

\textbf{Source qualifications.} The supplied record labels the problem open and reports that \(F\) has \(1\)-counter and \(3\)-counter graph-automatic structures, while ordinary automaticity and graph automaticity in the finite-state sense remain unresolved. These counter-automaton claims are supplied background, not assertions that the finite-state conditions above have been established. The record does not give the automaton conventions used for those counter structures. Its literature assessment and its formulation-verification labels are supplied curation, not independent verification in this draft.
\subsection{Short English statement}
Determine whether Thompson's group \(F\) is nonautomatic, as the source question anticipates, and whether its Cayley graph admits a presentation by a regular language with finite automata recognizing each generator relation. Ordinary automaticity requires representatives to be words in the generators; graph automaticity permits arbitrary word labels. The supplied counter-automaton results and open-status assessment require independent review.
\subsection{Sources}
\begin{itemize}
\item Local UnsolvedMath registry, permanent ID \(30003839\), exact source problem code \texttt{catalogue entry}, title: \emph{Graph-Automatic Structures for Thompson's Group \(F\)}. The supplied record belongs to the Oberwolfach Reports Open Problems collection and lists proposed year \(2018\); no individual proposer is supplied.
\item Supplied source citation: \emph{Cohomological and Metric Properties of Groups of Homeomorphisms of R} (2019), Oberwolfach Reports, pp.\ 1579--1633, DOI \texttt{10.4171/owr/2018/26}. \url{https://doi.org/10.4171/owr/2018/26}. The supplied background describes these reports as documenting mathematical workshops and their open-problem discussions.
\item Supplied literature source: Elder--Taback. \url{https://arxiv.org/abs/1501.04313}.
\item Supplied literature source: Taback--Younes. \url{https://arxiv.org/abs/1501.04315}.
\item Supplied literature source labelled ``Oberwolfach Report 2019''. \url{https://ems.press/content/serial-article-files/46748?nt=1}.
\item Literature checked at: \texttt{2026-08-21}. This exact date records supplied dated curation, not an independently verified assessment. The supplied assessment labels the problem open and states that \(F\) has \(1\)-counter and \(3\)-counter graph-automatic structures, whereas ordinary automaticity and finite-state graph automaticity remain unresolved. The background explicitly calls this dated literature triage derived from the supplied OWR curation and requests independent verification of time-sensitive status and cited claims.
\end{itemize}
\end{document}
